% =====================================================================
%  Chapter 1 — Introduction  (FULL PROSE)
%  Submission-ready written form of the section-level brief in
%  chapters/01_introduction.tex (kept unchanged) and ../draft/ch01_introduction.md.
%  Sources: thesis_writing_plan.md §3 (Ch.1), improvements.md, supervisor outline PDF,
%  MIPI I3C Basic v1.1.1. Ground-truth = current RTL/UVM (code wins over stale text):
%  realised flow is Xcelium + UVM 1.2 + SimVision; FPGA is future work only.
%
%  NOTE: not yet \include{}d in report.tex (the brief 01_introduction is). To compile
%  this version, swap that \include line to chapters/01_introduction_full.
% =====================================================================
\chapter{Introduction}
\label{ch:intro-full}

\section{I3C in SoC Trends and Motivation}
\label{sec:intro-motivation}

Modern systems-on-chip (SoCs) and embedded platforms aggregate an ever-growing number of sensors,
actuators, and housekeeping peripherals around a small set of low-pin-count serial buses. For more than
three decades the Inter-Integrated Circuit (\iic{}) bus has been the default interconnect for this role:
it needs only two wires, addresses multiple devices, and is backed by an enormous catalogue of
off-the-shelf parts. As device counts and sampling rates rise, however, its open-drain signalling caps
the practical clock rate and wastes static power, its manufacture-fixed addresses collide, and its lack
of a native interrupt forces wasteful polling --- limitations examined in
\cref{sec:bg-i2c}.

The MIPI Alliance introduced I3C (Improved Inter-Integrated Circuit) to lift these limitations while
preserving \iic{}'s two-wire simplicity and pin compatibility. The MIPI I3C Basic v1.1.1
specification~\cite{mipi_i3c_basic} --- the target standard of this thesis --- adds push-pull Single
Data Rate (SDR) signalling at up to \SI{12.5}{\mega\hertz}, run-time Dynamic Address Assignment (DAA),
and (in the full specification) In-Band Interrupts, while remaining backward compatible so that legacy
\iic{} Targets can share the same two wires as native I3C Targets.

These properties have made I3C a fast-growing choice for the sensor and peripheral fabric of contemporary
SoCs, and they make an I3C controller a compelling design exercise: it spans the full front-end
digital-IC pipeline, from reading a protocol specification, through register-transfer-level (RTL) design
of the protocol engine, to functional verification in a methodology-driven testbench. A
\emph{simplified, single-Controller} design is therefore an ideal, self-contained vehicle for studying
that pipeline end to end, which is the motivation for this thesis. \Cref{fig:soc-context} places such a
controller in its system context.

\begin{figure}[!htb]\centering
  \begin{tikzpicture}[
      font=\sffamily\small,
      block/.style={draw, thick, rounded corners, fill=blue!8,
                    minimum width=2.8cm, minimum height=1.2cm, align=center},
      target/.style={draw, thick, rounded corners,
                     minimum width=2.4cm, minimum height=1cm, align=center},
      wire/.style={thick},
      dot/.style={circle, fill=black, inner sep=1.2pt},
    ]
    % --- SoC interior: processor + controller ---
    \node[block] (cpu) at (0,0) {Application\\Processor};
    \node[block, fill=blue!18] (ctrl) at (5,0) {I3C Active\\Controller};
    \draw[wire] (cpu.east) -- (ctrl.west)
      node[midway, above=2pt, font=\scriptsize] {register bus};

    % --- SoC enclosure ---
    \node[draw, thick, dashed, rounded corners, fit=(cpu)(ctrl),
          inner sep=14pt] (soc) {};
    \node[anchor=north east, font=\itshape, inner sep=4pt] at (soc.north east) {SoC};

    % --- shared two-wire bus ---
    \coordinate (busL) at (-2.6,-3);
    \coordinate (busR) at (8.1,-3);
    \draw[wire] (busL) -- (busR);
    \node[anchor=south west, font=\scriptsize, inner sep=2pt] at (busL) {SCL / SDA};

    % --- controller drop to bus ---
    \draw[wire] (ctrl.south) -- (ctrl.south|-busL) coordinate (cdrop);
    \node[dot] at (cdrop) {};

    % --- target devices hanging off the bus ---
    \node[target, fill=green!12]  (t1) at (0,-4.7) {I3C\\Target 1};
    \node[target, fill=green!12]  (t2) at (3,-4.7) {I3C\\Target 2};
    \node[target, fill=orange!16] (t3) at (6,-4.7) {Legacy \iic{}\\Target};
    \foreach \t in {t1,t2,t3}{%
      \draw[wire] (\t.north) -- (\t.north|-busL) coordinate (\t d);
      \node[dot] at (\t d) {};
    }
  \end{tikzpicture}
  \caption{I3C in an SoC context: an application processor and an integrated I3C Active
    Controller, the controller driving native I3C Target devices and a legacy \iic{} Target
    device over a single shared two-wire (SCL/SDA) bus.}
  \label{fig:soc-context}
\end{figure}

\section{Problem Statement and Objectives}
\label{sec:intro-problem}

Although the I3C specification is freely available, learning the protocol \emph{from a production design}
is difficult. The open-source CHIPS Alliance \texttt{i3c-core}~\cite{chipsalliance_i3c} is a full-featured,
Host-Controller-Interface (HCI) compliant implementation integrated into the Caliptra root-of-trust
subsystem. Its production scope and platform integration make it a useful reference but obscure the
protocol path for study; in addition, several command-issue paths in its central finite-state machine
(FSM) remain as \texttt{TODO} stubs. Chapter~3 explains the qualitative design boundary derived from this
baseline, while Chapter~5 reports the quantitative comparison.

The objective of this thesis is therefore to derive, from that reference, a substantially simplified
single-Controller (Active Controller) design that (i)~retains the essential I3C SDR feature set,
(ii)~\emph{completes} the command FSM the reference leaves unfinished, and (iii)~is verified to work in a
modern, methodology-driven testbench. Concretely, the design targets the following deliverables: SDR
private write, read, and Immediate Data Transfer; \iic{} Fast-Mode legacy traffic on the same bus; ENTDAA
dynamic address assignment; the ENEC and DISEC Common Command Codes (CCCs); and a software-visible 32-bit
register interface --- all exercised by a reusable verification environment and a categorised regression
suite.

The methodology is a \emph{study-and-improve} flow rather than a clean-sheet design: the MIPI I3C Basic
specification and the reference architecture are studied first; the relevant subset is then re-implemented
and simplified in synthesisable SystemVerilog RTL; and the result is verified functionally in a Universal
Verification Methodology (UVM) 1.2 environment~\cite{uvm12} running on the Cadence Xcelium simulator, with
waveform debugging in SimVision. FPGA prototyping, listed as an optional step in the project outline, is
not undertaken in this work and is recorded only as future work (see the Conclusion). The resulting
controller and verification environment are evaluated against these objectives in \cref{ch:results}; this
introduction deliberately does not repeat the changing regression and implementation metrics reported
there. The precise scope of this work --- the functional requirements, the timing targets, and the
deliberate out-of-scope exclusions and their rationale --- is specified authoritatively in
\cref{ch:background-requirements} and is not duplicated here, so that later requirement changes cannot
produce a second, conflicting scope list.

\section{Contributions}
\label{sec:intro-contrib}

The contributions of this thesis are threefold. They are stated leading with the \emph{design} result
rather than a headline size figure; the quantitative code-size comparison with the reference --- and the
caveat that it follows largely from the exclusions defined in \cref{sec:req-oos} --- is reported once, in
\cref{ch:results}, where it can be presented with its full context.

\begin{enumerate}
  \item \textbf{A complete command FSM.} The flagship \texttt{flow_active} command FSM is implemented in
  full as a \FlowStateCount-state machine that drives SDR write, read, and immediate transfers, \iic{} legacy
  transfers, ENTDAA, and the abort path. This completes the command-issue paths that the
  CHIPS Alliance reference left as unimplemented \texttt{TODO} stubs (8 of its command-FSM states), so
  that every supported transaction can be followed end to end through the design.

  \item \textbf{A compact, studyable single-Controller design.} The controller spans the full stack ---
  from the physical (PHY) layer, through the protocol engine, to the register layer --- while retaining
  SDR, \iic{} Fast-Mode, and ENTDAA functionality. The reproducible size comparison is reserved for
  \cref{ch:results}.

  \item \textbf{A dual-agent UVM 1.2 verification environment.} A reusable testbench drives the design
  through two agents --- a register agent for the 32-bit host interface and an I3C Target agent for the
  bus --- and cross-checks behaviour with a scoreboard and SystemVerilog Assertions (SVA).

\end{enumerate}

\section{Thesis Organisation}
\label{sec:intro-organisation}

The remainder of this thesis is organised as follows. \Cref{ch:background-requirements} (Theoretical
Background and System Requirements) establishes the I3C and \iic{} protocol foundations and the UVM
methodology, then states the functional, performance, and verification requirements that bound the
design. \Cref{ch:architecture-rtl} (Architecture and RTL Design) is the flagship implementation chapter,
presenting the three-layer architecture and the RTL of each subsystem, culminating in the
\FlowStateCount-state \texttt{flow_active} command FSM. \Cref{ch:verif-method} (Verification Methodology
and UVM Environment) describes the directed-plus-coverage strategy and the dual-agent UVM environment
that exercises the design. \Cref{ch:results} (Results and Evaluation) reports the implementation metrics,
regression results, waveform evidence, and the comparison with the reference design, and
\cref{ch:conclusion} (Conclusion and Future Work) summarises the outcomes, the limitations, and the
future-work roadmap. Supporting reference material --- the register map, descriptor formats, FSM state
tables, the CCC table, regression logs, build-and-run instructions, and a glossary --- is collected in the
appendices.
